\ssr{ВВЕДЕНИЕ}

Графический планшет передаёт абсолютное положение стилуса, состояния наконечника и кнопок. Для управления интерфейсом по принципу мыши требуется преобразовать эти данные в относительные перемещения указателя и нажатия кнопок. Такое преобразование зависит от предыдущего отсчёта, присутствия пера в зоне считывания и выбранной чувствительности.

Актуальность задачи определяется необходимостью согласовать семантику периферийного устройства с интерфейсом прикладных программ. В Linux задачу можно решать пользовательской программой с uinput, HID-драйвером или собственным USB-драйвером. Штатные компоненты обеспечивают доставку и обработку ввода, однако наличие поддержки планшета само по себе не задаёт требуемый алгоритм относительного перемещения~\cite{kernelusb,uinput,libinput}.

Цель работы --- разработать загружаемый модуль ядра Linux, преобразующий ввод стилуса графического планшета в действия мыши.

Для достижения цели необходимо:
\begin{enumerate}
    \item сформулировать требования и сравнить способы обработки ввода;
    \item разработать функциональную модель, структуры и алгоритмы драйвера;
    \item реализовать приём USB-пакетов и передачу событий мыши;
    \item обеспечить согласованную остановку обработки и восстановление привязок драйверов;
    \item исследовать результаты работы и ограничения реализации.
\end{enumerate}

Объект работы --- обработка ввода периферийного устройства в Linux. Предмет работы --- преобразование отсчётов стилуса в относительные перемещения и состояния кнопок внутри модуля ядра. Аппаратная специализация определяется форматом отчёта и рассматривается в исследовательском разделе.
\clearpage
